\documentclass[11pt]{article}
\usepackage[T1]{fontenc}
\usepackage{lmodern}
\usepackage[margin=0.95in]{geometry}
\usepackage{amsmath,amssymb}
\usepackage{microtype}
\usepackage{enumitem}
\usepackage{xcolor}
\usepackage[colorlinks=true,linkcolor=blue!50!black]{hyperref}

\newcommand{\C}{\mathbb{C}}
\newcommand{\R}{\mathbb{R}}
\newcommand{\U}{U}
\newcommand{\SU}{\mathrm{SU}}
\newcommand{\abs}[1]{\lvert #1\rvert}
\newcommand{\norm}[1]{\lVert #1\rVert}
\newcommand{\Frob}{\mathrm{F}}
\newcommand{\key}[1]{\par\medskip\noindent{\normalsize\bfseries #1}\par\nobreak\vspace{-1pt}}
\setlist[enumerate]{leftmargin=1.5em,itemsep=3pt,topsep=3pt}

\title{\bfseries Step 3 of 5: certified cells inside $\U(2)$\\[5pt]
\large\normalfont Companion to \texttt{mub-outline.pdf}, Figure 1; preprint \S3.3--3.4 and \S4--5}
\author{}
\date{}

\begin{document}
\maketitle
\vspace{-2.6em}

\noindent
\textbf{What this step receives.} For each seed ball: a guarantee that every completion gives a point
$V\in\U(2)$ within $p\lesssim0.09$ of $\sqrt6$ times the compressed block, and nothing else.

\par\medskip\noindent\textbf{What it must produce.} A finite cover of $\U(2)$, and for each cell either a proof
that it holds no completion or a smaller certified region still containing all of them.

\key{1.\ Parametrise the group, then cut it}

Every $V\in\U(2)$ is $e(\alpha)$ times an element of $\SU(2)$, and every element of $\SU(2)$ is
$\bigl(\begin{smallmatrix}z&w\\-\bar w&\bar z\end{smallmatrix}\bigr)$ with $\abs z^2+\abs w^2=1$. Writing
$\abs z=\cos2\pi t$, $\abs w=\sin2\pi t$ and putting the two phases in $\beta,\gamma$ gives a
four-parameter chart
\[
V=e(\alpha)\begin{pmatrix} e(\beta)\cos(2\pi t) & e(\gamma)\sin(2\pi t)\\
-e(-\gamma)\sin(2\pi t) & e(-\beta)\cos(2\pi t)\end{pmatrix},
\qquad
0\le\alpha\le\tfrac12,\quad 0\le t\le\tfrac14,\quad 0\le\beta,\gamma\le1 .
\]
The chart's domain is a product of intervals, so it can be cut into closed \emph{cells} exactly as the
phase circle was cut into bins. A cell has a centre $U$ and a displacement $\theta$ in the four
angular coordinates, with $\abs{\theta_j}\le h=\pi/N_u$.

\key{2.\ Taylor with an honest remainder}

Two bounds make the grid usable, both uniform on the cell:
\[
\norm{V-U}_{\Frob}\le\sqrt6\,h,
\qquad
\norm{V-U-\mathrm dU[\theta]}_{\Frob}\le\frac{\sqrt{42}}{2}h^2 .
\]
The second is the useful one: inside a cell, the exact rotations are within $O(h^2)$ of a
\emph{linear} model, with a constant obtained by differentiating the four entries above.

\key{3.\ Certify, then inflate: the one pattern to take away}

Here is the pattern the whole proof runs on, and it is worth stating without symbols.

\emph{A check is designed for a point and is run on a model.} Suppose a test at the cell centre proves
that some quantity exceeds a threshold by a margin $\mu>0$. If the quantity moves by at most $c\,h^2$
across the cell, and $\mu>c\,h^2$, the test is true on the whole cell. So a test for exact points
becomes a test for regions, provided the remainder is bounded and the margin exceeds it. The computer
proposes the model; the inequality is what certifies; the fixed grid is what makes the remainder
uniform.

\key{4.\ Stage B: shrink the balls without losing a solution}

A cell cover is honest but coarse. Stage B refines it using the equations rather than the group:

\begin{enumerate}[label=(\alph*)]
\item at a retained ball's centre, compute numerical directions approximating the tangent space of
the Hadamard equations;
\item prove a \emph{residual identity} saying that every exact solution's displacement is within a
stated error of its projection onto those directions --- so the normal component is bounded by the
tangential coordinates;
\item cover the tangential coordinates by a grid, and inflate each cell by the normal bound.
\end{enumerate}

The result covers exactly the same exact solutions with smaller regions. Note what is \emph{not} done:
no manifold is parametrised, no component is classified, no dimension is assumed. A curved piece of
surface is covered by flat patches of known thickness, and that is all.

\key{5.\ Atlas and refinement}

Two bookkeeping steps complete the machinery. Section 4 collects certified phase balls from different
seeds into a single \emph{atlas}, so that a region already proved empty is not reproved; Section 5
handles the case where a ball is neither excluded nor resolved, by splitting it and reprocessing the
pieces, with a bound on how much the radii may grow in the process. Growth matters: if each split
inflated the region, the union could slowly stop containing the solutions it is supposed to contain.

\key{Where this step can fail, and why that is safe}

Refinement must terminate. If the grid is too coarse, a cell survives that should not; the fallback is
to split it again, and the proof is unfinished until the cell dies. If a cap on refinement is
exhausted, the region stays \emph{unresolved} --- a fact about the run, not about Hadamard matrices.
Nothing in Stages A or B can exclude a region it has not certified, so again the failure mode is
slow, not unsound.

\medskip
\noindent\textit{Note on sources.} The attached excerpt of the preprint ends in \S3.3, so the
parametrisation and the two Taylor bounds above are read from the paper, while the shape of Stage B and
of the atlas is taken from its \S1.2 summary. The reader with the full text should check \S4--5 for the
exact residual identity and the radius-growth statement.

\end{document}
